\documentclass[11pt]{article}

\usepackage[margin=25mm]{geometry}

\usepackage{amsmath,amssymb,fontspec,unicode-math}

\IfFontExistsTF{DejaVu Serif}{\setmainfont{DejaVu Serif}}{\setmainfont{Latin Modern Roman}}

\IfFontExistsTF{Latin Modern Math}{\setmathfont{Latin Modern Math}}{\setmathfont{DejaVu Math TeX Gyre}}

\usepackage{xurl,hyperref}

\setlength{\parindent}{0pt}\setlength{\parskip}{6pt}

\begin{document}

\title{Differentially Private Quantiles with Smaller Error}

\author{ELVIA preliminary AI analysis}\date{21 October 2026}\maketitle

Preliminary assessment: Minor.\par

Independent scientific review remains pending. Author adoption has not been requested.

Current integer counting noise fails the suffix-shift contract used to connect the implementation to the privacy proof. Historical indexing bugs are already fixed. Shared counting randomness and substitute-delta conversion are separate release/protocol issues; no final-output privacy violation is established.

\section*{Review scope}

Claim-focused reading of the original paper and relevant author-released passages; direct mathematical or static-source reasoning. No author program, earlier verifier, model or benchmark was executed during this check.

\section*{Problem 1: Signed integer tree noise has a forbidden suffix-shift parity}

Source: Lemma 3.1 Eq.(2), latest PDF p.5; current k\_ary\_tree\_noise.py, prefix\_noise\_terms

Assessment: supported; impact: minor

\subsection*{Source claim (paraphrase)}

The proof interface requires positive output probability to survive every unit suffix shift; the experimental implementation combines integer geometric-difference noise with signs.

\subsection*{Analysis}

At supported \(m=k=2\), height \(1\), independent integer nodes \(Z_0,Z_1,Z_2\) give \(N_1=Z_0+Z_1-Z_2\), \(N_2=Z_0+Z_1+Z_2\). Thus \(N_2-N_1=2Z_2\) is even. At ranks \((10000,20000)\), output equal to the ranks has positive probability because all nodes can be zero, while shifting the second output by \(1\) has zero probability. No finite exponential ratio can bound positive by zero. The same adjacent-difference parity appears in the height-one \(m=10\) experimental choice. Increased variance or widely separated ranks cannot repair support.

\subsection*{Finding and scope}

The current discrete module does not instantiate the required counting law. The continuous construction and existence lemma are not disproved; later randomized slice/quantile release may lose the internal parity information, so final-output nonprivacy is not established.

\subsection*{Proposed correction}

An ordinary dyadic-prefix mechanism with independently noised interval counts can satisfy the contract: a unit suffix shift changes one stream increment and at most \(T=\lceil\log_2m\rceil+1\) interval queries. Node mass proportional to \(e^{-\epsilon|z|/T}\) bounds the joint likelihood ratio by \(e^\epsilon\). This conservative alternative changes noise constants and needs full-pipeline integration/review.

A proposed correction is retained where stated; independent verification and author adoption remain unestablished.

\section*{Problem 2: The historical node-alias/prefix defect is already fixed}

Source: Conference supplement k\_ary\_tree\_noise.py; author commits of 10 September 2026; current prefix\_noise\_terms

Assessment: not\_issue; impact: minor

\subsection*{Source claim (paraphrase)}

The conference source reused sibling-index cache keys across parents and stopped at the first rather than last leaf boundary.

\subsection*{Analysis}

At \(m=8,k=2\), its coefficients force \(-N_1+2N_3-2N_4-N_5+2N_6=0\), whereas suffix shift \((0,1,1,1,1,1,1,1)\) changes that expression by \(1\). Current code uses parent-dependent absolute node keys and stops at rem equal to block size minus one, removing this rational support identity. These local fixes do not remove the separate integer parity above.

\subsection*{Finding and scope}

Known historical defect, cleared in the current source. Neither a historical benchmark change nor complete current privacy is inferred.

\section*{Problem 3: Nominal repetitions condition on one counting-noise draw}

Source: Section 6/Figure 1, latest PDF pp.9–10; current real\_data\_eps1.ipynb, cell 19; persistent tree cache

Assessment: supported; impact: minor

\subsection*{Source claim (paraphrase)}

One SliceQuantile object is constructed outside the \(200\)-repetition loop and its tree is not reset between queries.

\subsection*{Analysis}

Let \(Z\) be shared counting noise, \(Q_i\) random queries and \(E_i\) fresh single-quantile noise. Even if draws are conditionally independent, for loss \(Y_i=f(Q_i,Z,E_i)\), variance of the mean is \(\mathbb E[\operatorname{Var}(Y_i\mid Z)]/R+\operatorname{Var}(\mathbb E[Y_i\mid Z])\). Bootstrapping observations from one fixed \(Z\) cannot estimate the between-counting-noise term. New outer seeds do not clear an existing cache. Final exponential-mechanism outputs can still differ, so this is not a cached-final-answer claim.

\subsection*{Finding and scope}

The replication scheme supports conditional rather than fully randomized uncertainty unless the between-noise term vanishes. Numeric undercoverage or a ranking reversal is not established.

\subsection*{Proposed correction}

For independent complete-mechanism repetitions, create/reset the tree with a continuing nonreplayed RNG stream, while optionally pairing query seeds separately. State which sources of randomness the uncertainty covers.

A proposed correction is retained where stated; independent verification and author adoption remain unestablished.

\section*{Problem 4: Substitute delta is loosened instead of tightened}

Source: Theorem 4.2, latest PDF p.7; current slice\_quantile.py, is\_delta\_approximate\_DP

Assessment: supported; impact: minor

\subsection*{Source claim (paraphrase)}

The theorem gives final failure bound \(F\delta_{\rm int}\), \(F=1+e^{\epsilon_{\rm CC}+2\epsilon_{\rm EM}}\); source multiplies the requested delta by \(F\) before the internal bound.

\subsection*{Analysis}

At total epsilon \(1\), split \(0.5\), substitute adjacency, the correctly calibrated epsilon components are \(1/4,1/6\). Thus \(F=1+e^{7/12}=2.7920018257\). Requested \(10^{-16}\) needs internal \(10^{-16}/F=3.5816595491\times10^{-17}\), but source forwards \(F10^{-16}\), yielding only the conditional theorem bound \(F^2\,10^{-16}=7.7952741945\times10^{-16}\). No outer division in the inspected path cancels it.

\subsection*{Finding and scope}

The certification argument does not establish the requested substitute parameter. This is not a lower bound on actual final-output privacy loss, and the separate counting-contract issue remains.

\subsection*{Proposed correction}

Divide the requested final delta by the theorem’s factor once, keeping internal and target names distinct. This fixes accounting direction without proving privacy of the complete implementation.

A proposed correction is retained where stated; independent verification and author adoption remain unestablished.

\section*{Limits}

Preliminary AI assessment; independent human scientific adjudication remains pending.\par

Analytic witnesses and static release checks do not remeasure the published benchmark results.\par

\section*{Source identity}

arXiv:2505.13662v2, 12 December 2025

\url{https://arxiv.org/pdf/2505.13662v2}

Reviewed PDF SHA-256: \nolinkurl{8fb54de2548315cb3336c13d6c6bbd63a1f584463c7aae87957b963e8aba8b7f}

The SHA-256 identifies the exact retained PDF inspected in this rapid source/argument check. Original supplied report bytes are retained locally. The original PDF retrieval timestamp is taken from matching retained metadata.

\end{document}
